\chapter{Las cuentas}
\label{ch:ledgerch}

Este capítulo no narra. Los diez anteriores ya lo han hecho. Este audita: para cada
dominio, lo que produjeron sesenta y siete años, enunciado una sola vez, con sus
dos lados, y con la pregunta causal planteada en vez de dada por supuesta. Es la
bisagra del libro, porque es donde el registro se convierte en el conjunto de
restricciones dentro del cual tiene que diseñar la Parte III.

\section{Salud}

\emph{Lo que muestra el registro.} Cuba tomó un país con los mejores indicadores de
salud de América Latina y una mitad rural sin ninguno de ellos, cerró la brecha
rural-urbana y sostuvo los resultados nacionales durante sesenta años a través de
dos derrumbes económicos. La esperanza de vida pasó de unos 62 años en 1958 a los
setenta y tantos altos, convergiendo con Estados Unidos con una fracción del gasto.
La mortalidad infantil pasó de unos 37 por mil a un solo dígito y ahí se quedó
---incluido el período 1990--94, cuando siguió bajando mientras el producto
nacional caía un tercio---. Se eliminaron la polio, el sarampión, la difteria, la
rubéola y el tétanos neonatal. El modelo de prestación, un policlínico y después un
médico de familia responsables de una población definida, es el mecanismo, y otros
países lo han copiado deliberadamente porque funciona.

\emph{Lo que ha pasado desde 2019.} El sistema se está vaciando desde tres
direcciones a la vez. El suministro farmacéutico se ha derrumbado: la industria
nacional no consigue insumos y el Estado no puede pagar importaciones, y las
farmacias están vacías en proporciones que el propio ministerio de salud reconoce.
La planta física se ha degradado más allá del punto en que los indicadores de
resultado puedan sostenerse con el esfuerzo del personal ---hospitales sin agua,
electricidad ni esterilización fiables---. Y el personal se va, en la emigración
general y por las misiones médicas. Las cifras de médicos de Cuba siguen siendo
altas sobre el papel porque el papel cuenta a los médicos en misión en el
extranjero.

\emph{La lectura honesta.} El logro es real, no se compró con juegos de medición, y
es lo más valioso de este balance. Es también, en 2026, algo que se está perdiendo
en tiempo real, y no sobrevivirá otros cinco años en la trayectoria actual. Ese es
el argumento a favor del capítulo~\ref{ch:welfare} y la razón de que esté situado
temprano: el activo se deprecia más rápido de lo que puede escribirse el plan.

\section{Educación}

\emph{Lo que muestra el registro.} Analfabetismo eliminado en un año y mantenido
eliminado. Escolarización primaria y secundaria universal. Rendimiento medido por un
organismo externo dos desviaciones estándar por encima de la media latinoamericana,
con el cuartil cubano más pobre por encima del promedio regional ---que es la cifra
que importa, porque dice que el sistema comprime en vez de reproducir la ventaja de
origen---. Enseñanza superior gratuita y ampliada en un orden de magnitud.
Rendimiento femenino por delante del masculino desde los años ochenta.

\emph{El lado del costo, y la erosión.} El currículo fue un sistema de formación
política, campos enteros estuvieron cerrados o reducidos a catecismo, y Cuba tiene
en consecuencia ingenieros y médicos excelentes y muy poca gente formada para pensar
cómo funciona un mercado, un tribunal o una empresa ---lo que es una restricción
directa sobre la aplicabilidad de la Parte III, y no una cuestión de gusto---. Y el
sistema se erosiona ahora: los maestros se van al sector privado y al extranjero, la
matrícula en las escuelas pedagógicas ha caído, y ha aparecido una economía informal
de repaso privado para llenar el hueco, que es la primera etapa clásica de un sistema
de dos niveles.

\emph{La lectura honesta.} El mejor activo que tiene Cuba, aquello con lo que están
hechos tanto la segunda vía como el sector biotecnológico de la Parte III, y aquel
cuya depreciación es menos visible año a año, porque el declive de un sistema
escolar aparece en los resultados una década después.

\section{Alimentos}

\emph{Lo que muestra el registro.} Un país con buen suelo, agua suficiente, ciclo
de cultivo todo el año y una profesión agronómica formada importa entre dos tercios
y cuatro quintas partes de lo que come, y no puede pagarlo.\unv{} Una parte
sustancial de la tierra cultivable está ociosa o bajo marabú. La industria
azucarera que era tres cuartas partes de las exportaciones tan recientemente como
en 1989 produjo zafras de cientos de miles de toneladas a mediados de la década de
2020 y el país importa azúcar.

\emph{Las causas, por orden de magnitud.} Monopolio estatal de acopio que fija
precios por debajo del costo; tenencia que no puede venderse, hipotecarse ni
heredarse con fiabilidad; sin mercado de insumos; sin transporte ni cadena de frío
funcionales; y una moneda en la que nadie quiere que le paguen. Todas ellas son
domésticas y legales.

\emph{Y el dinero fue a otra parte.} En 2024, la construcción vinculada al turismo
se llevó el 37,4 por ciento de la inversión nacional y la agricultura el 2,7 por
ciento ---una razón de catorce a uno, en un país que pedía leche para niños al
Programa Mundial de Alimentos---. Sea lo que sea que se discuta sobre la política
económica cubana, esta asignación se eligió.

\emph{La lectura honesta.} Este es el caso más claro del libro en que el embargo
explica muy poco. Cuba puede comprar alimentos a quien quiera en el mundo y se los
compra a Estados Unidos al contado. La razón por la que no cultiva los suyos es el
diseño del sistema agrícola doméstico, que el capítulo~\ref{ch:land} desarma.

\section{Energía}

\emph{Lo que muestra el registro.} Sesenta años de generación eléctrica a partir de
hidrocarburos importados en condiciones políticas ---soviéticas, después
venezolanas--- quemados en grandes plantas térmicas centrales, la mayoría de las
cuales ha superado ya su vida de diseño. Cuando el suministro político flaqueó en
1990 el país se apagó; cuando volvió a flaquear después de 2019 el país se apagó, y
en octubre de 2024 la red falló como sistema, a escala nacional, por primera vez.

\emph{La lectura honesta.} La energía es donde convergen las tres dependencias
externas de Cuba, y es la restricción vinculante sobre todo lo demás ---agua,
refrigeración, procesamiento de alimentos, salud, turismo, industria,
conectividad---. Es además el único dominio donde la tecnología ha cambiado lo
suficiente desde 2010 como para que la respuesta disponible hoy no sea la que
estaba disponible entonces. El capítulo~\ref{ch:energy} va en consecuencia por
delante de todos los capítulos sectoriales de la Parte III.

\section{Demografía}

\emph{Lo que muestra el registro.} La fecundidad cayó por debajo del reemplazo en
1978 y nunca ha vuelto; a mediados de la década de 2020 rondaba 1,2 a 1,4, de las
más bajas del mundo.\unv{} Más de una cuarta parte de la población tiene más de
sesenta años. Y entre uno y dos millones de personas, desproporcionadamente en edad
de trabajar y de tener hijos, se fueron entre 2021 y 2025 de una población de once
millones.

\emph{La lectura honesta.} Es la restricción más dura del libro y la menos
reversible. No hay cohorte en camino, de modo que ninguna política adoptada en 2026
cambia la población en edad de trabajar antes de la década de 2040 salvo por
migración. Toda propuesta de la Parte III está por tanto limitada por la mano de
obra y no por el capital, y un sistema de pensiones con tasas de reemplazo europeas
no es financiable con esta razón de dependencia a ninguna tasa impositiva plausible
---razón por la cual el capítulo~\ref{ch:welfare} tiene que decirlo en vez de darlo
por resuelto---.

\section{Derechos}

\emph{Lo que muestra el registro, sin inflarlo y sin excusarlo.} Ninguna prensa
independiente desde 1961. Ninguna organización política legal distinta de una.
Ningún sindicato independiente. Ningún tribunal que haya fallado jamás contra el
ejecutivo en materia política. Ejecuciones sumarias en 1959 sin conteo publicado.
Campos de trabajo para creyentes, homosexuales e inconformes, 1965--68. Una purga
cultural por reasignación administrativa, 1971--76. La emigración castigada con
confiscación y trabajo forzado durante dos décadas, y turbas organizadas contra los
vecinos que se iban en 1980. Setenta y cinco personas condenadas a hasta veintiocho
años en 2003 y tres fusiladas el mismo mes. Treinta y siete personas ahogadas frente
a La Habana en 1994 sin investigación independiente. Cientos de procesados tras el
11 de julio de 2021, con condenas de cuatro a treinta años, menores incluidos.

\emph{Lo que también debe decirse.} La escala no es la de las peores dictaduras de
la región, e inflarla es un flaco favor: Cuba no operó escuadrones de la muerte, no
desapareció gente por decenas de miles, y la cifra de veinte mil muertos atribuida a
Batista es ella misma una fabricación que se lleva reciclando sesenta años. El
crimen violento es bajo, la policía no mata gente a tasas latinoamericanas, y los
cubanos se mueven de noche por sus propias ciudades de un modo que los ciudadanos de
la mayor parte del hemisferio no pueden. Sesenta años de partido único produjeron un
Estado represivo, no uno asesino, y la distinción importa cuando el
capítulo~\ref{ch:polity} diseñe la transición, porque el arreglo que un país
necesita después de Argentina no es el que necesita después de esto.

\section{La cuestión del embargo}

Es la pregunta a través de la cual se discute todo lo demás, y la respuesta honesta
tiene dos mitades que suelen publicarse por separado.

\emph{Lo que el embargo sí cuesta.} No es nada, y la versión desdeñosa es
equivocada. Transporte marítimo: un buque que hace escala en un puerto cubano no
puede entrar en un puerto de Estados Unidos durante 180 días, lo que saca a Cuba de
las rutas caribeñas eficientes y encarece el flete de todo. Finanzas: la aplicación
extraterritorial ---de la que las multas multimillonarias impuestas a bancos
europeos son la demostración--- hace que los bancos internacionales rechacen el
negocio cubano por política antes que por ley, de modo que una transacción cubana es
cara, lenta o imposible incluso donde es legal. Bienes médicos e industriales:
cualquier producto con más de una proporción mínima de contenido norteamericano
queda atrapado, lo que en la práctica cubre buena parte del equipamiento médico
moderno. Crédito: Cuba está excluida del Fondo Monetario Internacional, el Banco
Mundial y el Banco Interamericano de Desarrollo, lo que no es solo una pérdida de
préstamos sino de la asistencia técnica y del sello de aprobación que otros países
usan para atraer capital privado. Y la Helms-Burton convirtió todo eso de
instrumento ejecutivo en estatuto, de modo que un presidente no puede negociarlo.

\emph{Lo que el embargo no explica.} Cuba comercia libremente con el resto del
mundo. La Unión Europea, China, Canadá, Brasil, Rusia, México, España, Vietnam y
Venezuela no han embargado nada, y Cuba les ha comprado a todos. La restricción
vinculante sobre las importaciones cubanas no es el permiso; es que Cuba no puede
pagar, y no puede pagar porque está en impago desde 1986, congeló las cuentas de
empresas extranjeras en 2009, y tiene atrasos con casi todo proveedor que le haya
dado crédito. Eso es un historial crediticio y no tiene autor extranjero. Tampoco
explica el embargo las prohibiciones domésticas: la abolición del pequeño negocio en
1968, la negativa de personalidad jurídica a las empresas privadas hasta 2021, la
ausencia de mercado mayorista, el monopolio estatal de acopio en la agricultura, la
prohibición a los cubanos de entrar en hoteles hasta 2008, el permiso de salida
hasta 2013, la concesión de licencias por lista enumerada, o la agencia empleadora
que se queda con el diferencial cambiario de los trabajadores de las empresas
mixtas. Todas ellas se escribieron en La Habana, y todas pueden derogarse en La
Habana.

\emph{Cómo sostener ambas.} La formulación que usa este libro es que el embargo fija
el costo del fracaso cubano y el sistema cubano fija su probabilidad. Una economía
cubana bien diseñada bajo embargo sería más pobre que la misma economía sin él
---apreciablemente más pobre, y la diferencia se mide en puntos porcentuales de
crecimiento al año, no en los totales acumulados fantásticos que publica cada uno de
los dos bandos---. Pero no sería una economía incapaz de iluminar sus ciudades o de
alimentarse, porque varios países de la región logran ambas cosas con peores
dotaciones y sin asistencia externa alguna. La razón de que la Parte III de este
libro se concentre en las instituciones domésticas no es que el embargo carezca de
importancia. Es que el embargo es la única variable que los cubanos no controlan, y
un diseño que espera por ella no es un diseño.

\section{El patrón}

A lo largo de sesenta y siete años la misma secuencia ocurre cuatro veces, y
nombrarla es el hallazgo contra el que hay que construir la Parte III.

\emph{La presión produce la reforma.} 1980, 1993--95, 2008--2013, 2021: en cada
caso una emergencia material obligó al Estado a legalizar algo que había prohibido
---mercados campesinos, dólares y trabajo por cuenta propia, tierra en usufructo y
venta de casas, empresas privadas con personalidad jurídica---.

\emph{La reforma funciona.} Aparece la comida; se estabiliza la moneda; crece el
sector privado; se surten las tiendas. En todos los casos el resultado medido fue
positivo y visible en cuestión de meses.

\emph{La presión cede y la reforma se retira.} 1986, 1996--2004, 2016--17,
2024--25: mercados cerrados, licencias congeladas, impuestos subidos, categorías
eliminadas, topes de precio impuestos, y los reformadores criticados por
concentración de riqueza.

\emph{La retirada llega justo antes del momento en que es menos asequible.} La
Rectificación se lanzó en 1986 y la Unión Soviética se derrumbó en 1990. La
congelación de licencias llegó en agosto de 2017, catorce meses antes del reapriete
norteamericano. El ataque a los nuevos importadores privados empezó en 2024, en
medio de la peor escasez desde 1993.

La razón del patrón no es económica y fingir lo contrario hace perder el tiempo a
todo el mundo. Cada reforma crea gente cuyo ingreso no viene del Estado, y eso es
precisamente lo que el sistema está organizado para impedir. Las reformas se
presentan por tanto siempre como concesiones, siempre revocables, y siempre se
revocan ---y como se sabe que son revocables, nadie invierte detrás de ellas, lo que
garantiza que rindan poco, lo que suministra la evidencia para revocarlas---.

Ese es el hallazgo, y se convierte directamente en el requisito de diseño con el que
abre la Parte III: \emph{el problema vinculante en Cuba no es qué políticas adoptar
sino cómo volverlas irrevocables}. Cada medida que este libro propone ya se ha
intentado en Cuba de alguna forma. La mayoría funcionó. Todas se revirtieron. Un
diseño cuyo único contenido sea una lista mejor de políticas no ha entendido los
últimos treinta años, y el capítulo~\ref{ch:polity} está situado en segundo lugar en
la Parte III por esa razón: sin un mecanismo creíble de compromiso, nada de lo que
venga después vale la pena escribirlo.

\begin{ledger}{1959--2026}
\emph{Logrado.} Un sistema de salud que cerró la brecha rural-urbana, convergió con
resultados de país rico con dinero de país pobre, y sostuvo sus indicadores a través
de una contracción del 35 por ciento. Un sistema educativo cuyos resultados medidos
superan a su región en dos desviaciones estándar y que comprime en vez de reproducir
la ventaja de origen. Analfabetismo terminado en un año. Desegregación racial legal y
aplicada. Seguridad social universal, propiedad de vivienda casi universal, y una de
las tasas de crimen violento más bajas del hemisferio. Un sector biotecnológico que
desarrolló y desplegó sus propias vacunas cuando países mucho más ricos no pudieron.
Una contribución decisiva al fin del dominio sudafricano en Namibia. Y un Estado que
nunca ha perdido la capacidad de llegar a cada cuadra de su territorio, que es la
razón de que absorbiera 1993 sin hambruna.

\emph{Costo.} Una canasta exportadora sin cambio estructural entre 1958 y 1989 y
destruida desde entonces. Tres patrones externos sucesivos y tres fracasos idénticos
en diversificar bajo ellos. Azúcar por debajo de la producción del siglo XIX y
alimentos importados en dos tercios del consumo en un país fértil. Una red eléctrica
que falló como sistema en 2024. Un impago sin resolver desde 1986. Seis décadas sin
una prensa, un sindicato, un partido o un tribunal que no fueran del Estado. Campos,
purgas, el exilio como castigo, turbas organizadas contra los vecinos, y condenas de
treinta años en 2021 por caminar por la calle. Fecundidad por debajo del reemplazo
desde 1978 y entre uno y dos millones de personas idas en cinco años. Y un historial
de reformas en el que todo lo que funcionó se retiró después, que es la razón de que
nada de ello se acumulara.
\end{ledger}
